
\documentclass[12pt]{article}
\usepackage{amsmath, amsthm, amsfonts, mathtools, xcolor, caption}

\usepackage{minimalJ}

\numberwithin{equation}{section}

\usepackage[style=numeric,sorting=none,giveninits=true]{biblatex}
\addbibresource{bibliography.bib}

\usepackage{etoolbox}
\usepackage[margin=1in]{geometry}

\DeclareMathAlphabet{\mymathbb}{U}{BOONDOX-ds}{m}{n}

%\usepackage[colorlinks, allcolors=blue, hypertexnames=false]{hyperref}
%\usepackage[noabbrev, nameinlink]{cleveref}

\usepackage{comment}

%%%%%%%%%%%%%%%%%%%%%%%%%%%%%%%%%%%%%%%%%
%       Theorem Environments            %
%%%%%%%%%%%%%%%%%%%%%%%%%%%%%%%%%%%%%%%%%

\newtheorem{theorem}{Theorem}[section]
\newtheorem{lemma}[theorem]{Lemma}
\newtheorem{corollary}[theorem]{Corollary}
\newtheorem{proposition}[theorem]{Proposition}
\theoremstyle{definition}
\newtheorem{remark}{Remark}
\newtheorem{definition}{Definition}
\newtheorem{notation}{Notation}
\newtheorem{example}{Example}
\newtheorem{conjecture}[theorem]{Conjecture}
\newtheorem{openproblem}[theorem]{Open Problem}

%%%%%%%%%%%%%%%%%%%%%%%%%%%%%%%%%%%%%%%%%
%               New Commands            %
%%%%%%%%%%%%%%%%%%%%%%%%%%%%%%%%%%%%%%%%%

\newcommand{\ev}[1]{( #1 )} 
\newcommand{\ew}[1]{\left( #1 \right)} 

\newcommand{\N}{\mathbb{N}}
\newcommand{\Z}{\mathbb{Z}}
\newcommand{\C}{\mathbb{C}}
\newcommand{\Q}{\mathbb{Q}}
\newcommand{\R}{\mathbb{R}}



%%%%%%%%%%%%%%%%%%%%%%%%%%%%%%%%%%%%%%%%%
%               Document                %
%%%%%%%%%%%%%%%%%%%%%%%%%%%%%%%%%%%%%%%%%

\let\underbrace\LaTeXunderbrace
\let\overbrace\LaTeXoverbrace


\begin{document}

\section*{Challenge 5: Erdős--Hajnal Conjecture}
\refstepcounter{section}



\begin{definition}[Cliques and Independent sets]
Given a simple graph $G$, a set of vertices of $G$ is a \textit{clique} if its vertices are pairwise adjacent. It is an \textit{independent set} if its vertices are pairwise non-adjacent. 
\end{definition}

\begin{definition}[$H$-free graphs]
Given graphs $G$ and $H$, we say that $G$ is \textit{$H$-free} if it does not contain an induced subgraph isomorphic to $H$. 
\end{definition}

\begin{conjecture}[Erdős–Hajnal Conjecture \mbox{\cite{erdos_hajnal}}]
For each graph $H$, there is a constant $c_H > 0$ such that every $H$-free graph $G$ with $n$ vertices has a clique or independent set of size at least $n^{c_H}$. 
\end{conjecture}



Given $r\in \mathbb{N}$, we denote the path with $r$ vertices by $P_r$.
We are interested in the following parametrized version of the conjecture, where $H=P_r$ is a path. 
\begin{conjecture}[Erdős–Hajnal Conjecture for paths]
Let $r\in \mathbb{N}$. There is some $c_r > 0$ such that every $P_r$-free graph $G$ with $n$ vertices has a clique or independent set of size $n^{c_r}$.
\end{conjecture}

The conjecture has been proved for $r=5$ by Nguyen, Tung, Scott and Seymour in \cite{nguyen2026induced}


%%%%%%%%%%%%%%%%%%%%%%%%%%%%%%%%%%%%%%%%%%%%%%%%%%
\printbibliography
%%%%%%%%%%%%%%%%%%%%%%%%%%%%%%%%%%%%%%%%%%%%%%%%%%
\end{document}
